\scanpage{26}
\begin{center}- 64 -\end{center}

\uline{Définition 2.} \uline{Les hypothèses et notations étant celles de la prop.1}, $B$ \uline{s'appelle la fermeture intégrale de} $A$ \uline{dans} $R$; \uline{si} $A=B$, $A$ \uline{est dit intégralement fermé dans} $R$; \uline{si} $B=R$, $R$ \uline{est dit entier sur} $A$.
\uline{Si} $A$ \uline{est anneau d'intégrité}, \uline{la fermeture intégrale} $C$ \uline{de} $A$ \uline{dans son corps des fractions s'appelle la clôture intégrale de} $A$; \uline{si} $A=C$, $A$ \uline{est dit intégralement clos}.

D'après la prop.1, pour qu'un anneau d'intégrité soit intégralement clos, il faut et il suffit qu'il soit intersection d'anneaux de valuation de son corps des fractions.

\uline{Remarque}: Une intersection d'anneaux de valuation d'un corps $K$ n'est pas nécessairement un anneau ayant $K$ pour corps des fractions. Par exemple, les anneaux de valuation de $k(X)$ qui contiennent le corps $k$ ont une intersection réduite à $k$ (exemple 3 du § 3,n°1).
\margin{$S$}

\uline{Exemples}. [ici, prière de transcrire les exemples 1) et 2) de la rédaction précédente, p. \unsure{32|38}--39].

\uline{Proposition 2} (``transitivité''). \uline{Soient} $A,B$ \uline{deux sous-anneaux d'un anneau} $R$, \uline{tel que tout élément de} $B$ \uline{soit entier sur} $A$. \uline{Alors tout élément de} $R$ \uline{entier sur} $B$ \uline{est entier sur} $A$.

C'est immédiat en vertu de l'équivalence de a) et \unsure{e)|c)} du th.1.

\uline{Proposition 3.} \uline{Soient} $R,B,A$ \uline{trois anneaux tels que} $R\supset B\supset A$; \uline{soit} $S$ \uline{une partie non vide de} $A$, \uline{multiplicativement stable}, \uline{ne contenant pas} $0$. \uline{Alors}, \uline{si} $B$ \uline{est entier sur} $A$, $B_S$ \uline{est entier sur} $A_S$; \uline{si} $B$ \uline{est la fermeture intégrale de} $A$ \uline{dans} $R$, $B_S$ \uline{est celle de} $A_S$ \uline{dans} $R_S$.

Rappelons, pour justifier cet énoncé, qu'on a convenu d'identifier $A_S$, $B_S$ avec des sous-anneaux de $R_S$ (v. la remarque qui suit le coroll.1 de la prop.2, § 1); si $\varphi$ est l'homomorphisme canonique de $R$ dans $R_S$, on a $A_S=\varphi(A)\varphi(S)^{-1}$, $B_S=\varphi(B)\varphi(S)^{-1}$. Supposons $B$ entier sur $A$;
\margin{25}
